\section{问题重述}

\subsection{问题背景}

来袭导弹沿指向假目标的方向匀速飞行，无人机则在接令后选定固定航向和速度，并通过控制烟幕弹的投放与延时起爆，在导弹与真目标之间形成半径为 $10\,\mathrm{m}$ 的有效烟幕球。问题的核心并非烟幕球是否靠近真目标，而是烟幕球能否在有限寿命内切断导弹到真目标的视线。

\subsection{问题提出}

题目给出了三枚来袭导弹、五架无人机和真假目标的初始位置，以及各类飞行器的运动参数与烟幕云团的有效范围和持续时间。本文需要由这些离散的初始数据建立统一的运动学与视线遮蔽关系，并在不同装备规模和任务对象下，依次完成遮蔽时长计算、投放参数优化、协同策略设计与资源分配。具体问题可归纳如下。

\textbf{问题一：}研究单架无人机投放单枚烟幕弹干扰单枚导弹的基准情形。需要综合描述导弹、无人机、烟幕弹和烟幕云团的运动过程，建立烟幕云团与导弹视线之间的几何判据，并据此计算给定投放策略产生的有效遮蔽时长。

\textbf{问题二：}将问题一的给定参数计算推广为单弹投放优化问题。需要把无人机航向、飞行速度、投放时刻和起爆延迟纳入决策，在满足运动学与烟幕有效性约束的条件下，以延长有效遮蔽时间为目标，确定单枚烟幕弹的最优投放方案。

\textbf{问题三：}进一步考虑单架无人机连续投放多枚烟幕弹的情形。除沿用单弹遮蔽模型外，还需处理无人机航向和速度一致、相邻投放至少间隔 $1\,\mathrm{s}$ 等时序约束，协调三枚烟幕弹的投放与起爆时刻，使多个遮蔽区间形成尽可能长的联合有效时间。

\textbf{问题四：}将干扰平台由单架无人机扩展为三架无人机，并由各机投放一枚烟幕弹共同干扰 M1。需要利用不同初始位置带来的空间互补性，协调三架无人机的飞行与投放参数，形成面向同一导弹的协同遮蔽策略。

\textbf{问题五：}面向五架无人机、三枚来袭导弹的多平台多目标场景，需要同时解决无人机与导弹之间的任务分配，以及各架无人机所携烟幕弹的数量、投放时机和空间位置安排。在单机最多投放三枚烟幕弹等约束下，建立资源调度与时空协同模型，获得整体有效遮蔽效果尽可能优的投放方案。

\section{问题分析}

\subsection{问题一的分析}

问题一的输入参数均由题面给定，不含待优化变量。其计算链由导弹轨迹、无人机轨迹、烟幕弹自由落体、烟幕云团下沉和视线遮蔽五部分组成。本文保留真目标半径 $7\,\mathrm{m}$、高度 $10\,\mathrm{m}$ 的圆柱几何，要求烟幕遮住完整目标在导弹视场中的投影，并以代表点口径作为对照。由于遮蔽状态在阈值附近发生切换，采用事件根求解进入与离开时刻，可避免固定时间步长造成的系统误差。
\subsection{问题二的分析}

问题二沿用问题一已经验证的运动学方程与完整圆柱遮蔽判据，但将给定策略改为四维连续决策。无人机一旦选定航向和速度便不能再次调整，因此投放点与起爆点并非相互独立，而是共同受同一条水平航迹、投放时刻和引信延时约束。为减少冗余变量，本文以航向角、飞行速度、起爆时刻和引信延时编码策略，再反算投放时刻及两处空间坐标。目标函数是阈值条件 $D(t)\le 10\,\mathrm{m}$ 所对应时间集合的测度，具有非光滑、多峰和大面积零值等特点；故采用带未命中距离引导的多初值差分进化，并对不同航向扇区中的可行盆地独立细化和交叉比较，最后通过加密采样、事件求根及边界扰动复核候选策略。

\subsection{问题三的分析}

问题三在问题二单弹优化的基础上增加了三弹时序耦合。三枚烟幕弹共享 FY1 的航向和速度，但各自具有投放时刻与引信延时；相邻两次投放至少间隔 $1\,\mathrm{s}$，使各弹的空间位置、起爆时刻和有效区间不能独立选取。评价指标也不能把三枚弹的单独遮蔽时长直接相加，否则重叠区间会被重复计数。本文分别求出各弹满足完整圆柱遮蔽判据的时间集合，再以三者并集的测度作为目标，并用松弛量编码投放间隔以保证候选解天然满足时序约束。针对非光滑目标和投放下界附近的窄可行盆地，采用多盆地差分进化、时序边界构造和高密度事件求根联合求解；当多个参数组合达到相同联合时长时，再比较有效区间重叠裕量和任务完成时刻，选取更稳健、执行更紧凑的代表方案。

\subsection{问题四的分析}

问题四把单机多弹改为三机各一弹。三架无人机的初始位置差异较大，可到达的有利遮蔽空域和时间窗也随之不同，因此不能简单复制问题二的 FY1 参数。本文先分别搜索 FY1、FY2、FY3 对 M1 的单弹有效窗口，再以三个窗口并集的测度作为联合目标。若窗口互不相交，则联合时长等于各单弹时长之和；若出现重叠，则进一步联合调整投放和起爆时刻，使不同平台分别承担早、中、晚阶段的遮蔽任务。该分解既保留完整圆柱几何判据，又把十二维联合问题化为三个四维子问题和一个时间窗协调问题。

\subsection{问题五的分析}

问题五同时包含离散任务分配和连续轨迹优化。每架无人机一旦确定航向和速度，其全部烟幕弹必须共享同一条航迹；每架最多投放三枚，相邻投放至少间隔 $1\,\mathrm{s}$；同一枚烟幕弹只计入一个被指派导弹的遮蔽收益。为避免直接在高维混合空间中盲目搜索，本文先对全部无人机--导弹组合进行单弹效能探针，筛除难以形成完整遮蔽的低收益组合；随后按边际收益构造任务分配，并在每个导弹分组内优化共享航迹和投放时序。总目标取三枚导弹各自有效时间并集测度之和，这既避免同一导弹上的重叠重复计时，也能反映多目标场景下的总体保护水平。

\endinput

\section{模板的基本使用}

\subsection{LaTeX 发行版和编辑器安装}

要用 \LaTeX{} 来完成建模论文，首先要确保正确安装一个 \LaTeX{} 的发行版本。

\begin{itemize}
    \item Mac 下可用 Mac\TeX{}
    \item Linux 下可用 \TeX{}Live ；
    \item windows 下可用 \TeX{}Live 或 Mik\TeX{} ;
\end{itemize}

具体安装可以参考 \href{https://github.com/OsbertWang/install-latex-guide-zh-cn/releases/}{Install-LaTeX-Guide-zh-cn} 或者其它靠谱的文章。另外可以安装一个易用的编辑器，例如 \href{https://mirrors.tuna.tsinghua.edu.cn/github-release/texstudio-org/texstudio/LatestRelease/}{\TeX{}studio} 。

使用该模板前，请阅读模板的使用说明文档。

\subsection{文档类选项说明}

根据要求，电子版论文提交时需去掉封面和编号页。可以加上 \verb|withoutpreface|  选项来实现，即：
\begin{tcode}
    \documentclass[withoutpreface]{cumcmthesis}
\end{tcode}
这样就能实现了。打印的时候有超链接的地方不需要彩色，可以加上 \verb|bwprint| 选项。

在撰写论文的过程中可加上 \verb|draft| 选项以选择不嵌入图片来提高编译的速度，最终提交的电子版文档请去掉 \verb|draft| 选项以保证图片和代码被嵌入！若选择使用 \verb|draft| 选项请务必保证检查最终文档的图片和代码是否被嵌入，避免因图片和代码未被嵌入而影响参赛成绩。

\subsection{目录说明}

另外目录也是不需要的，将 \verb|\tableofcontents| 注释或删除，目录就不会出现了。

\subsection{编译方式}

编译记得用 \verb|xelatex|，而不是用 \verb|pdflatex|。在命令行编译的可以按如下方式编译：
\begin{tcode}
  xelatex main
\end{tcode}
或用 \verb|latexmk| 来编译，更推荐这种方式。
\begin{tcode}
  latexmk main
\end{tcode}

下面给出写作与排版上的一些建议。
